% !TEX root = root.tex

{\color{red}(Here goes all notation that needs to be explained.)}
In this paper, we consider the setting of real time model learning over
networked  systems. The system is modeled in the form
\begin{equation}
        \mathcal{P}:\, y[k] = \phi[k]\theta_{\text{true}} + w[k]
        \label{eq:sysmodel}
\end{equation}
{\color{red}(Bracket notation for time indices $k$ is maybe a bit old fashioned.
Later on when I start involving attack signals more explicitly, I will however be using so many
subindices that I'd like to save the space.)}
where $y[k]$ are the measured outputs, $\phi[k]$ is a row vector of
the regressors consisting of current and past data, $w[k] \in \mathbb{R}^n$
are i.i.d. zero-mean gaussian measurement errors with covariance $\gamma^2$, and
$\theta_{\text{true}} \in \mathbb{R}^n$ are the true system parameters.
{\color{red}(In the case of an ARX model, it might be more prudent to say that
$w[k]$ is a disturbance, and thus show up in $\phi[k]$?)}
The system (\ref{eq:sysmodel}) can be modeled using any model structure that
is linear in the parameters. In this paper, we will mainly be considering
ARX and FIR model structures. In the ARX case, we have
\begin{equation}
        \begin{split}
                \phi[k] &= \begin{bmatrix}
                        y_{[k-1:k-n]} & u_{[k-1:k-n]}  
                \end{bmatrix} \\
                y_{[k-1:k-n]} &= \begin{bmatrix}
                        y[k-1] & y[k-2] & \cdots & y[k-n]
                \end{bmatrix} \\
                u_{[k-1:k-n]}  &= \begin{bmatrix}
                        u[k-1] & u[k-2] & \cdots & u[k-n]
                \end{bmatrix}
        \end{split}
\end{equation}
and similarly in the FIR case we have
\begin{equation}
        \phi[k] = \begin{bmatrix}
                u_{[k-1:k-n]}  
        \end{bmatrix}
\end{equation}
where $n$ is the order of the system. \cite{ljung_system_1999}

The physical plant $\mathcal{P}$ is connected to the operator over network.
Inputs are sent from the network to the plant, where the outputs are
measured and then sent back to the operator, again over the network.
These signals are then used to estimate a model of the system.
It is assumed in this paper that at any time, the network might be
compromised by a \textit{strategic adversary}\cite{chong_tutorial_2019}.
The adversary might have access to the input/output signals to learn about
the system, and the ability to corrupt the signals transmitted over network.
Thus, the inputs received
at the plant side $\tilde{u}[k]$ and the output measurements received at the
operator side $\tilde{y}[k]$ might not be the true signals, and we have
\begin{align}
        \tilde{y}[k] = y[k] + \delta_y[k] \\
        \tilde{u}[k] = u[k] + \delta_u[k]
\end{align}
as illustrated in figure \ref{fig:adversary-schematic}.

\begin{figure}
        \centering
        % !TEX root = root.tex

\usetikzlibrary{shapes.symbols}

\tikzstyle{block} = [draw, fill=white, rectangle, 
    minimum height=1cm, minimum width=1cm]
\tikzstyle{sum} = [draw, fill=white, circle, node distance=1cm]
\tikzstyle{branch} = [draw, fill=black, circle,
minimum width=1pt, node distance=2cm, scale=0.5]
\tikzstyle{input} = [coordinate]
\tikzstyle{output} = [coordinate]
\tikzstyle{pinstyle} = [pin edge={to-,thin,black}]
\tikzstyle{network} = [cloud, cloud puffs=7, draw, aspect=2, node distance=3cm]

\begin{tikzpicture}[auto, node distance=2cm,>=latex]
    \node [block, name=P] {\large$\mathcal{P}$};
    \node [name=adversary, below of=P]
        {\includegraphics[width=1cm]{figs/angryfaic.png}};
    % \draw(0,-2.7) node{Adversary};
    \node [network, left of=adversary, name=networku] {Network};
    \node [network, right of=adversary, name=networky] {Network};
    % \node [block, below of=adversary, name=C] {\large $\mathcal{C}$};
    \node [block, below of=P, name=E, node distance=5cm] {\large Estimator};
    \node [branch, left of=E,node distance=6cm, name=ubranch] {};
    \node [coordinate, below of=ubranch](u){};

    %Adversary lines
    \draw[->] (adversary) to [bend right]
        node[pos=0.3]{$\delta_u[k]$} (networku);
    \draw[->] (networku) to[bend right]
        node[pos=0.3]{$u[k]$} (adversary);

    \draw[->] (adversary) to [bend right]
        node[pos=0.3]{$\delta_y[k]$} (networky);
    \draw[->] (networky) to[bend right]
        node[pos=0.3]{$y[k]$} (adversary);

    %C-E line
    % \draw[->] (E) -- node{$\hat{\theta}$} (C);

    %Block lines
    \draw[->] (P) -| node{$y[k]$} (networky);
    % \draw[->] (networky) |- node[pos=0.49]{$\tilde{x}_k$} (C);
    \draw[->] (networky) |- node{$\tilde{y}[k]$} (E);
    % \draw[->] (C) -| node {$u_k$} (networku);
    \draw[<->] (networku) |- (E);
    \draw[->] (u) -- node{$u[k]$} (ubranch);
    \draw[->] (networku) |- node {$\tilde{u}[k]$} (P);
\end{tikzpicture}


% \node [input, name=input] {};
% \node [block, right of=input] (Cr) {$C_r(s)$};
% \node [sum, right of=Cr, node distance=1.5cm] (esum) {};
% \node [block, right of=esum] (controller) {C(s)};
% \node [block, right of=controller, 
%         node distance=2cm] (system) {G(s)};
% \node [sum, right of=system, node distance=1.5cm] (dsum) {};
% \node [block, above of=dsum,
%     pin={[pinstyle]above:d}] (Gd) {$G_d(s)$};
% \draw [->] (Cr) -- (esum);

% \draw [->] (Gd) -- (dsum);

% \draw [->] (controller) -- node[name=u] {$u$} (system);
% \node [output, right of=dsum] (output) {};
% \node [input, below of=u] (measurements) {};

% \draw [draw,->] (input) -- node {$r$} (Cr);
% \draw [->] (esum) -- node {$e$} (controller);
% \draw [->] (system) -- (dsum);
% \draw [->] (dsum) -- node [name=y] {$y$}(output);
% \draw [->] (y) |- (measurements) -| (esum);
% \draw [->] (measurements) -| node[pos=0.99] {$-$} 
% node [near end] {$y_m$} (esum);
        \caption{Block schematic of the networked adaptive control system, with strategic adversary.}
        \label{fig:adversary-schematic}
\end{figure}